\section{Evidence Lanes For Intervention Thinking}

This chapter uses the book's evidence policy explicitly.

\begin{itemize}[nosep]
  \item \textbf{Strong lane.} State-transition evidence from consciousness and
        neural-dynamics work supports the general claim that systems can occupy
        distinct global regimes and move between them through structured
        transitions. This supports treating intervention as regime-shaping
        rather than mere self-talk.
  \item \textbf{Medium lane.} Metastability, attractor, integrator, causal
        emergence, and nonlinear-control perspectives support the design idea
        that repeated local inputs can stabilize or redirect future behaviour.
  \item \textbf{Speculative lane.} Frontier consciousness-physics proposals may
        inspire hypotheses about large-scale organization, but they do not
        justify practical intervention rules in this chapter.
  \item \textbf{Design inference.} The concrete protocols below are applied
        synthesis. They are informed by the strong and medium lanes, but they
        remain design constructions rather than direct experimental theorems.
\end{itemize}

\begin{proposition}[Interventions should target the transition geometry]
When a failure is recurrent, the first design question should usually be:
which alteration most cheaply changes the probability of entering the target
state during a real decision window?
\end{proposition}

\begin{discussion}
This proposition is stronger than ``make the task easier.'' It says the
intervention should be evaluated in transition terms. If the current state is a
locally stable basin, then a useful intervention either lowers the barrier,
changes the timing of attempted movement, or alters the local environment so
that the old state becomes less self-restoring. That logic is consistent with
the book's state-transition framing and with the broader dynamical-systems
language of attractors, metastability, and bifurcation-like regime shifts.
\end{discussion}
